%!TEX program = uplatex
\RequirePackage{plautopatch}
\documentclass[uplatex,dvipdfmx]{jlreq}
\usepackage{amsmath,amssymb,amsfonts,amsthm}
\usepackage{enumerate}

\pagestyle{empty}

\begin{document}

%======================================================================
% 講演1：関口次郎「ベキ零軌道の魅力」
%======================================================================
\begin{center}
{\large\bfseries ベキ零軌道の魅力}\\[6pt]
\hspace*{\fill}{\large\bfseries 関口 次郎（東京農工大学 工学部）}
\end{center}

単純リー環にはベキ零元というやっかいな代物が住んでいます。
ベキ零元がなければ、その構造ももう少しあつかいやすものに
なっていたでしょうが、存在するために複雑になりました。
逆に言えば、
そのために構造が豊富になりました。

${\mathbf sl}_2$ はもっとも簡単な単純リー環です。
一般の単純リー環 ${\mathbf g}$ に対して、
${\mathbf sl}_2$ から ${\mathbf g}$ へのリー準同型の共役類は
1950年代に Dynkin が分類しました。
この Dynkin の論文~[2] で今日 Dynkin 図形と呼ばれる図形を使って、
単純リー環の場合の分類を記述しました。
さらに、Kostant がベキ零軌道と ${\mathbf sl}_2$ から ${\mathbf g}$
へのリー準同型
の共役類との関係を明確にしたことで (cf.~[4])、
複素数体上のベキ零軌道の分類が完成しました。

有理2重点と Dynkin 図形の関係が注目されるようになり、
それを単純リー環のベキ零元との関係へと深化させたのが
Grothendieck と Brieskorn の成果です (cf.~[6])。
一方では、Harish-Chandra は不変固有超関数の研究において、
Kostant の構造論を使って、ベキ零多様体に台をもつ
不変固有超関数は存在しないという、
Local integrability theorem の鍵となる結果を導いています
(cf.~[3])。

この講演では複素数体上の単純リー環のベキ零元について、
基本的な事柄と上述のことを解説する予定です。
([1] にはベキ零軌道についての基本的な事柄がまとめられています。)

\begin{center}\textsc{参考文献}\end{center}
\begin{enumerate}[{[1]}]
  \item D.H.~Collingwood and W.M.~McGovern:
    \textit{Nilpotent Orbits in Semisimple Lie Algebras},
    Van Nostrand Reinhold Mathematics Series.
  \item E.~Dynkin:
    Semisimple subalgebras of semisimple Lie algebras,
    \textit{Amer.\ Math.\ Soc.\ Transl.\ Ser.~2}, \textbf{6} (1957), 111--245.
  \item Harish-Chandra:
    Invariant distributions on Lie algebras,
    \textit{Amer.\ J.\ Math.}, \textbf{86} (1964), 271--309.
  \item B.~Kostant:
    The principal three-dimensional subgroup and
    the Betti numbers of a complex simple Lie group,
    \textit{Amer.\ J.\ Math.}, \textbf{81} (1959), 973--1032.
  \item 草場公邦:
    行列特論, 基礎数学選書21（裳華房）.
  \item 松澤淳一:
    特異点とルート系, すうがくの風景6（朝倉書店）.
\end{enumerate}

\newpage

%======================================================================
% 講演2：中島啓「1. 超ケーラー多様体入門 / 2. インスタントンと巾零軌道」
%======================================================================
\begin{center}
{\large\bfseries 1. 超ケーラー多様体入門\\
2. インスタントンと巾零軌道}\\[6pt]
\hspace*{\fill}{\large\bfseries 中島 啓（京都大学 大学院理学研究科）}
\end{center}

Kronheimer は、Nahm 方程式の解、あるいは
${\mathbb R}^4$ 上の $SU(2)$-不変
なインスタントン解のモジュライ空間が、複素リー環の巾零軌道（一般には、
その Slodowy slice）になることを証明し、特に後者が超ケーラー多様体の構造
を持つことを証明した。超ケーラー多様体は、微分幾何学の片隅で研究されて
いる分野であるが、このように数学の中心分野で研究されている対象が超ケー
ラー多様体になることが、ときどき起こる。他の有名な例が、非可換ホッジ理
論における、複素射影多様体上の半単純局所系のモジュライ空間である。
（Hitchin、藤木による。）

さて、超ケーラー多様体の理論の応用として、小林--ヒッチン型の定理と組み
合わせると、二つの一見異る複素多様体が、実は同じ超ケーラー多様体から来
ていることが分かることがある。Vergne は、このアイデアに基づき、
Kostant--関口対応を解釈した。

本講演では、超ケーラー多様体入門~[1, 2]（とはいっても入門以
上の理論は存在しない）から始めて、
Kronheimer~[3] や Vergne~[4]
の結果を紹介する予定である。

\begin{center}\textsc{参考文献}\end{center}
\begin{enumerate}[{[1]}]
  \item N.J.~Hitchin:
    Monopoles, minimal surfaces and algebraic curves,
    \textit{S\'eminaire de Math\'ematiques Sup\'erieures} 105,
    Les Presses de l'Universit\'e de Montr\'eal, 1987.
  \item N.J.~Hitchin, A.~Karlhede, U.~Lindstr\"om and M.~Ro\v{c}ek:
    Hyperk\"ahler metrics and supersymmetry,
    \textit{Comm.\ Math.\ Phys.}, \textbf{108} (1987), 535--589.
  \item P.B.~Kronheimer:
    Instantons and geometry of the nilpotent variety,
    \textit{J.\ Differential Geom.}, \textbf{32} (1990), 473--490.
  \item M.~Vergne:
    Instantons et correspondence de Kostant--Sekiguchi,
    \textit{C.R.\ Acad.\ Sci.\ Paris S\'er.\ I Math.}, \textbf{320} (1995), 901--906.
\end{enumerate}

\newpage

%======================================================================
% 講演3：落合啓之「ベキ零軌道の関口対応」
%======================================================================
\begin{center}
{\large\bfseries ベキ零軌道の関口対応}\\[6pt]
\hspace*{\fill}{\large\bfseries 落合 啓之（名古屋大学 大学院多元数理科学研究科）}
\end{center}

舞台は、
実半単純リー環とその involution（$=$位数2の自己同型）である。
半単純対称空間の接空間が典型的な例である。
このとき2種類のベキ零軌道の間に全単射の対応が存在することを、
関口次郎の原論文に沿って紹介する。

ここでの議論の鍵となるのは、
ベキ零元から $\mathfrak{sl}(2)$-triple を構成する
Jacobson--Morozov の定理、
つまりリー環の構造論である。
この構成の特徴を述べて、
次の中島氏の講演で紹介される幾何学的構成との
長所短所を比較する。

舞台が特別な場合、すなわち、
複素半単純リー環とその実形を指定する involution の場合には、
関口対応は
「Kostant--Sekiguchi 対応」と呼ばれることもある。
半単純リー群の無限次元表現論の文脈で良く用いられるのはこの場合である。

また、関連して、松木対応
（旗多様体上のある群の作用に関する軌道分解と
別の群に関する軌道分解が1対1に対応する）
についても述べておく予定である。

\begin{center}\textsc{参考文献}\end{center}
\begin{enumerate}[{[1]}]
  \item D.~Barbasch and M.R.~Sepanski:
    Closure ordering and the Kostant--Sekiguchi correspondence,
    \textit{Proc.\ Amer.\ Math.\ Soc.}, \textbf{126} (1998), 311--317.
  \item T.~Matsuki:
    The orbits of affine symmetric spaces under the action
    of minimal parabolic subgroups,
    \textit{J.\ Math.\ Soc.\ Japan}, \textbf{31} (1979), 331--357.
  \item J.~Sekiguchi:
    Remarks on nilpotent orbits of a symmetric pair,
    \textit{J.\ Math.\ Soc.\ Japan}, \textbf{39} (1987), 127--138.
\end{enumerate}

\newpage

%======================================================================
% 講演4：竹内潔「特性サイクルの理論とその実リー群の無限次元表現論への応用について1・2」
%======================================================================
\begin{center}
{\large\bfseries 特性サイクルの理論とその実リー群の\\
無限次元表現論への応用について1・2}\\[6pt]
\hspace*{\fill}{\large\bfseries 竹内 潔（筑波大学 数学系）}
\end{center}

「複素半単純リー環の表現の指標公式は Kazhdan--Lusztig 多項式
の特殊値で表現されるであろう」という Kazhdan--Lusztig 予想は、
1980年代初頭に Brylinski--Kashiwara および Beilinson--Bernstein により、
リー環の表現と旗多様体上の $D$-加群を結びつける
画期的な手法で解決された（[7] 参照）。

その後、1980年代後半に入って、Kazhdan--Lusztig 予想
の解決で使われた手法を用いて、実リー群の無限次元表現を
同変構成可能層を用いて幾何学的に構成するプログラムが
柏原により提案された。これは Borel--Weil によるコンパクト
群の表現の幾何学的構成を、非コンパクト群の表現の
場合に一気に拡張するものである。

本講演では、柏原による一連の予想の解決の過程において
Schmid--Vilonen~[4, 5, 6] により解かれた、「表現（の
Harish-Chandra 加群）の特性多様体と表現の指標超関数の波面集合が
Kostant--Sekiguchi 対応で対応するであろう」という
Barbasch--Vogan 予想~[1] の解決のあらましについて、紹介する。

証明には、柏原により導入された構成可能層の特性サイクル
と呼ばれる概念が本質的に用いられる（[2] 参照）。
Mirkovic--Uzawa--Vilonen~[3]
による層の松木対応の超局所版を確立して、モーメント写像でリー環の
上に落とすことで、懸案であった対応が得られる。すなわち、
無限次元表現の代数的不変量（特性多様体）と解析的不変量
（指標超関数の波面集合）が、幾何学的な対象（特性サイクル）
を仲立ちにして結びつけられる、という世にも美しい理論について
紹介したい。

\begin{center}\textsc{参考文献}\end{center}
\begin{enumerate}[{[1]}]
  \item D.~Barbasch and D.A.~Vogan~Jr.:
    The local structures of characters,
    \textit{J.\ Funct.\ Anal.}, \textbf{37} (1980), 27--55.
  \item M.~Kashiwara and P.~Schapira:
    \textit{Sheaves on Manifolds},
    Grundlehren der Math.\ Wiss.\ 292, Springer-Verlag, 1990.
  \item I.~Mirkovic, T.~Uzawa, and K.~Vilonen:
    Matsuki correspondence for sheaves,
    \textit{Invent.\ Math.}, \textbf{109} (1992), 231--245.
  \item W.~Schmid and K.~Vilonen:
    Characteristic cycles of constructible sheaves,
    \textit{Invent.\ Math.}, \textbf{124} (1996), 451--502.
  \item W.~Schmid and K.~Vilonen:
    Two geometric character formulas for reductive Lie groups,
    \textit{Jour.\ A.M.S.}, \textbf{11} (1998), 799--867.
  \item W.~Schmid and K.~Vilonen:
    Characteristic cycles and wave front cycles of representations of reductive Lie groups,
    \textit{Ann.\ Math.\ (2)}, \textbf{151} (2000), 1071--1118.
  \item 谷崎俊之、堀田良之：$D$-加群と表現論, シュプリンガー現代数学シリーズ.
\end{enumerate}

\vfill

\end{document}
